% Article VI. Owner: indice_articulos. Self-contained expository addition.
% Provenance: RESULTADO_RECUPERADO; new documentary checks kept separate.
% REV2/fuente/modulos/15_comparacion_cuantitativa_integral.md (complete).
% REV2/fuente/modulos/16_precision_robustez_y_contraste.md (complete).
% REV2/fuente/modulos/19_realizacion_exhaustiva_estado_por_estado.md (complete).
% Integral: integracion_83/masas_p0/corpus/source/masas/
% 19_comparacion_cuantitativa_posterior.tex, portion preceding \endinput.
% Integral: integracion_83/masas_p0/tex/apendice_inventario_471_384_sucesor.tex.
% Depends on the construction of the atlas, signature, towers and material realisation
% incorporated by the editor in the preceding sections. Does not replace it.

\section{Realisation of observables and metrological comparison}
\label{sec:vi-contraste-metrologico}

A mass value acquires its meaning within the structure that
produces it and the representation in which it is measured. The
APP--TRIT--TPK construction first determines admissible trajectories, their oriented
records and the signature readers; the towers and material realisation
act on these objects. Experimental comparison takes place at the end of
this composition. Its function is to compare observables of the same
type, preserve the identity of the state and quantify agreement in a
common chart. This position makes it possible to study together the mathematical
structure, the evaluated coordinate and its physical interpretation.

This section brings together three levels with different functions:
the internal domain of routes and their quotients; the classes and individualised
evaluations in the mass source owners; and the external inventory of
observables used as a reference. The order of exposition preserves the
constructive direction of the article. External quantities appear as
terms of comparison, while internal values preserve the signatures,
operations and conditions fixing their evaluation.

\subsection{Structural domain, material classes and external inventory}
\label{subsec:vi-dominios-contraste}

Let \(\mathcal C\) be the set of APP cells and \(\mathcal R\) its domain
of oriented routes. The identification of a route contains its cell
coordinates and the two initial orientations:
\[
 \gamma=(r,c,h_0,v_0),\qquad
 r,c\in\{1,\ldots,9\},\quad h_0,v_0\in\{-1,+1\}.
\]
Hence, \(\#\mathcal C=81\) and \(\#\mathcal R=81\cdot4=324\).
The \(56\) trajectory collections and \(13\) multisections belong
to quotients and structural groupings of the same domain, with the
relations and readers established earlier. These four cardinalities
describe different internal objects. Assigning a physical
representation may preserve an entire fibre; a species label does not
by itself impose a pointwise choice within it.

The public table of \(23\) classes organises realisations by sector.
Its function is to relate internal residence to representation,
charge, spin, colour, flavour and observable type. A class may
comprise several species, and a catalogue identity may bring together
different charge states. The correspondence is expressed through
typed maps and preserves these multiplicities.

The frozen metrological inventory contains \(471\) mass records and
\(384\) width records. The corpus documents \(1170\) identities in \(438\)
groups, of which \(403\) contain at least one of these observables.
The sum
\[
 471+384=855
\]
counts external summary-table observables. The \(81/56/13/324\) atlas,
the \(23\) classes and the \(855\) records thus belong to three levels
of description. Their association preserves the type of each set, rather
than converting their cardinalities into a single number of particles or predictions.

\begin{center}
\begin{tabular}{p{0.49\linewidth}rrr}
\toprule
Catalogue collection & Groups & Masses & Widths\\
\midrule
General & 46 & 53 & 9\\
Quarks & 6 & 8 & 1\\
Mesons & 201 & 243 & 225\\
Baryons & 149 & 166 & 149\\
External gravitational record & 1 & 1 & 0\\
\midrule
Total with observable & 403 & 471 & 384\\
\bottomrule
\end{tabular}
\end{center}

The gravitational row preserves the identity appearing in the external source;
its presence in the inventory does not define an HMT elementary particle.
The remaining \(35\) groups preserve their catalogue identity, although
the snapshot assigns them no summary-table mass or width. Documentary
integrity encompasses both facts: the recorded quantities and the identities
for which that snapshot does not express one of those quantities.

\subsection{The observable before the numerical coordinate}
\label{subsec:vi-tipos-observable}

A material output is specified through the state, representation,
operator and reading procedure. For a fibre \(F_X\), let
\(\widehat M_X\) be the mass operator already constructed, and let \(P_{X,a}\)
be the projector corresponding to the preparation state or channel \(a\).
For \(\Psi_X\) normalised in its image, the spectral distribution is
\[
 \mu_{X,a}(B)
 =\langle\Psi_X,
   P_{X,a}E_{\widehat M_X}(B)P_{X,a}\Psi_X\rangle,
\]
where \(E_{\widehat M_X}\) is the operator's spectral measure. This
expression preserves the prepared state and the multiplicity of the fibre.
A single scalar coordinate for every state prepared in that subspace
is obtained under the condition
\[
 \operatorname{Ran}(P_{X,a})\subset\operatorname{Dom}(\widehat M_X),
 \qquad \widehat M_XP_{X,a}=m_XP_{X,a}.
\]
Indeed, every vector in that subspace is then an eigenvector of the
mass operator with the same value \(m_X\), and its spectral measure is
concentrated at that value. If the compression equality
\(P_{X,a}\widehat M_XP_{X,a}=m_XP_{X,a}\) is used, one additionally requires
that \(P_{X,a}\) reduce \(\widehat M_X\); in finite dimension, this is equivalent
to commutation. Without that condition, the compression may fix only
the mean value. When the spectral measure retains several values, the
output is presented as a distribution or multiplet with its weights.

For a fixed density \(\rho\) in a finite fibre, let \(m_j\)
be the distinct eigenvalues and \(E_j\) their spectral projectors.
The weights \(p_j=\operatorname{tr}(\rho E_j)\geq0\), with
\(\sum_jp_j=1\), determine
\[
 \overline m=\sum_jp_jm_j,\qquad
 \operatorname{Var}_{\mu}(m)=\sum_jp_j(m_j-\overline m)^2.
\]
\paragraph{Scalar-concentration criterion for the fixed state.}
The variance vanishes if and only if all values with positive weight
coincide. Indeed, every summand is non-negative; the sum is zero
precisely when \(m_j=\overline m\) for every \(j\) with
\(p_j>0\). This determines when a mass distribution admits
a scalar representation without losing the dispersion of the prepared state.
The criterion concerns the fixed density and differs from the
operator condition on the entire subspace. The electronic case uses its
operator on the central fibre, as constructed earlier;
the rank-one property of an arbitrary projector is insufficient to concentrate
the spectral measure of another operator.

Mass and width use coordinated readers of different types. For a
resonant state, a pole convention may be expressed in energy
coordinates as
\[
 z_X=M_Xc_{\rm int}^{2}-\frac{i}{2}\Gamma_X.
\]
The real part is the mass--energy coordinate; \(\Gamma_X\) has
the dimension of energy and measures the width. Once the time chart is fixed,
\(\tau_X=\hbar_{\rm int}/\Gamma_X\). The relation assumes the pole
and survival regime for which it was constructed; a rate, a width and
a lifetime retain their units and conversion rules. This
distinction explains why the inventory preserves two observable classes.

The same principle determines comparison by sector. Running quark masses
are expressed with a scheme \(\mathscr S\) and a scale \(\mu\);
resonances distinguish pole mass, Breit--Wigner parametrisation and
on-shell convention. Neutrinos distinguish mass eigenvalues,
mass-squared differences, flavours and mixing parameters. One-sided
bounds preserve the documented confidence level and statistical convention.
An upper bound is not converted into a positive central value,
nor a difference of squared masses into an absolute mass.

\subsection{Fixing the construction and matching records}
\label{subsec:vi-congelacion}

The realisation of a species \(X\) transmits the necessary information
from its internal antecedents. For comparison purposes, we may collect it
in the record
\[
 \mathfrak D_X=
 \bigl(F_X,\rho_X,P_{X,a},\mathcal L_X,
       \sigma_X,\eta_X,\widehat M_X,\mathcal U_X\bigr),
\]
where \(\rho_X\) identifies the representation, \(\mathcal L_X\) the
genealogical record, \(\sigma_X\in\mathbb Z^5\) the signature,
\(\eta_X\) the tower section and \(\mathcal U_X\) the unit
chart. All these data receive their meaning from the
APP--TRIT--TPK operations and the realisations proved in the article. Collecting them
in a tuple is a documentary operation: their constructions and
proofs remain in the sections that produce them.

The external reference is represented separately by
\[
 o=(\mathrm{id},\mathrm{species},\mathrm{observable},
      v,\mathrm{unit},\mathscr S,\mu,
      \mathrm{uncertainty},\mathrm{limit},\mathrm{confidence},
      \mathrm{source},\mathrm{edition}).
\]
The association \((\mathfrak D_X,o)\) must be established through
physical identity and observable type. In an extensive table, this
operation uses identity and provenance keys; numerical proximity
does not serve as a key. The record preserves the
particle or antiparticle label, the group possessing the observable and
the identity of the measured property.

The independence of the selection procedure is examined by hiding or
permuting the metrological column and checking that \(\mathfrak D_X\)
remains unchanged. This test is meaningful for a generator whose
effective composition can be executed; merely preserving an output
file does not by itself constitute such a test. The cryptographic hash
fixes the artefact being compared and makes modifications detectable, but
causal independence belongs to the operation and its dependencies.

Statuses are assigned to the domain and operation supporting each
assertion. A row may simultaneously document an internal operator,
an evaluation given a signature and an external comparison, without identifying
these three levels. Every subsequent realisation is linked to the
definition, proposition or table establishing its scope. In particular, the
catalogue's documentary classifications apply to the snapshot specified by
each record and do not change the mathematical status of a subsequent
construction.

\subsection{Metrological correspondence of the twenty-three classes}
\label{subsec:vi-clases-metrologia}

The following table preserves the source's \(23\) classes and specifies the
form of reading required by each upon reaching the metrological interface. Its
function differs from the operator classification in
Proposition~\ref{vi:prop:producto-fibrado-especie-ruta} and the documentary inventory
in Tables~\ref{tab:vi-estados-clases} and~\ref{tab:vi-clases}:
here the comparable quantity, realisation convention and
accompanying data are distinguished. Flavour and sector names
guide the correspondence; the transmitted mathematical object remains in
the corresponding fibres, representations and readers.

\begingroup
\small
\setlength{\tabcolsep}{3pt}
\begin{table}[!htbp]
\centering\fontsize{9.2}{10.6}\selectfont\renewcommand{\arraystretch}{1.04}
\caption{Metrological reading of the twenty-three classes and type of transmitted object.}
\label{tab:vi-clases-metrologia}
\begin{tabular}{p{0.20\linewidth}p{0.25\linewidth}p{0.47\linewidth}}
\toprule
Class & Representatives & Object documented for the reading\\
\midrule

Charged lepton 1 & Electron and positron & Central route, conjugation and beta coordinate.\\
Charged lepton 2 & Muon and antimuon & Fibre, readers and evaluation of the declared signature.\\
Charged lepton 3 & Tau and antitau & Fibre, readers and evaluation of the declared signature.\\
Neutrino 1 & Electron flavour and conjugate & Neutral fibre and mass--flavour realisation.\\
Neutrino 2 & Muon flavour and conjugate & Neutral fibre and mass--flavour realisation.\\
Neutrino 3 & Tau flavour and conjugate & Neutral fibre and mass--flavour realisation.\\
Up-type quark 1 & \(u,\bar u\) & Flavour and colour fibre; scheme and scale chart.\\
Down-type quark 1 & \(d,\bar d\) & Flavour and colour fibre; scheme and scale chart.\\
Up-type quark 2 & \(c,\bar c\) & Flavour and colour fibre; scheme and scale chart.\\
Down-type quark 2 & \(s,\bar s\) & Flavour and colour fibre; scheme and scale chart.\\
Up-type quark 3 & \(t,\bar t\) & Flavour and colour fibre; scheme and scale chart.\\
Down-type quark 3 & \(b,\bar b\) & Flavour and colour fibre; scheme and scale chart.\\
Electromagnetic gauge & Photon & Null chart and helicity representation.\\
Chromatic gauge & Gluons & Null chart and adjoint colour representation.\\
Charged gauge & \(W^+,W^-\) & Route, signature and pole convention of the realisation.\\
Neutral gauge & \(Z^0\) & Route, neutral mixing and pole convention.\\
Scalar & Higgs & Scalar representation, signature and pole convention.\\
Mesonic composite & Mesons and conjugates & Constituents, binding record and composite reading.\\
Baryonic composite & Baryons and conjugates & Trilinear composition, binding and composite reading.\\
Fibonacci topological & \(1,\tau\) & Fusion ring and Frobenius--Perron dimension.\\
Ising topological & \(1,\psi,\sigma\) & Fusion rules and sectoral representation.\\
Cyclic topological & \(\mathbb Z/6\mathbb Z\) & Winding characters and braid representation.\\
Virtual & Intermediate sections & Finite horizon, prolongation and obstruction.\\
\bottomrule
\end{tabular}
\end{table}
\endgroup

For charged leptons, the reading leads to the mass of the central section
or to evaluation of the selected fibre. For neutrinos, it also preserves
the transformation between mass and flavour and the mass-squared differences; it therefore
does not reduce to an isolated scalar. For quarks, every figure is accompanied by the
renormalisation scheme and scale. The electromagnetic and
chromatic classes respectively express an experimental bound and the null gauge
section; the \(W\), \(Z\) and scalar classes require the declared pole
or resonance convention. Composites preserve mass, width and constituent
identity. The three topological classes are compared through their
fusion, braiding and gap laws once the dynamical realisation is fixed;
their dimensions and characters belong to their own codomains. The virtual
class preserves the pole, obstruction and prolongation horizon. This
reading determines which object reaches metrology without conflating the
twenty-three classes with twenty-three universal scalar masses.

\subsection{Metrological snapshot, units and transcription accuracy}
\label{subsec:vi-corte-metrologico}

The snapshot identified by the corpus as Particle Data Group
2026 is used, with a cutoff date of \(15\) January \(2026\) and a documented
consultation on \(26\) July \(2026\). This section reproduces that
documentary state: it does not incorporate a subsequent catalogue update.
The references in the evaluation supplement and those in the historical grid
retain their individual provenance; they are not automatically assigned the edition
of the general inventory.

Each value is preserved as its original decimal string, together with its
presentation, unit, uncertainty and provenance key. This decision
prevents an intermediate floating-point conversion from changing the last
digits or removing significant zeros from the transcription. The original
string and the typographical representation are distinct data: the
former reproduces the file's arithmetic exactly; the
latter communicates the result with precision appropriate for reading.
The digital representation digits of a reference do not increase its
experimental precision.

Conversion between MeV and GeV uses the exact factor \(10^3\), applied
to both the value and its uncertainty. To present a mass, the unit is written
MeV/\(c^2\) or GeV/\(c^2\); to present rest energy,
MeV or GeV. The relation \(E=mc^2\) is the chart allowing passage from one
reading to the other. Comparing coordinates requires fixing that chart in
both columns.

The inventory snapshot expressly records the scheme and scale of its external
records as not supplied. A property name
may contain an indication such as “Breit--Wigner”; it is preserved literally,
but that partial indication does not complete the metadata required by the comparison.
Empty cells, unavailable-data markers and
not-applicable markers remain distinct. A documentary unknown
is not replaced by zero.

\subsection{Individualised evaluations and reading conditions}
\label{subsec:vi-evaluaciones-documentadas}

Table~\ref{tab:vi-evaluaciones-veinte} brings together \(20\) numerical
presentations: the electronic beta coordinate, two null outputs and
\(17\) evaluations conditional on declared descriptors. Among
the latter, \(11\) correspond to elementary sectors and \(6\) to
composites. This count expresses the content of that documentary presentation;
the fibres, collections and records of the domain retain their own extent.

The table preserves all twenty entries. To facilitate reading,
four modalities are identified: \(\mathrm{B}\), electronic beta
coordinate; \(\mathrm{N}\), null chart; \(\mathrm{F}\), evaluation
conditional on a complete signature; \(\mathrm{A}\), conditional angular
evaluation. The subscript \(c\) denotes material composition. The
condition of each entry belongs to its result type and is also preserved
in the electronic catalogue accompanying the text.

\begingroup
\small
\setlength{\tabcolsep}{3pt}
\begin{table}[!htbp]
\centering\fontsize{8.8}{10.2}\selectfont
\setlength{\tabcolsep}{4pt}\renewcommand{\arraystretch}{1.04}
\caption{Twenty individualised evaluations, with their chart and comparison modality.}
\label{tab:vi-evaluaciones-veinte}
\begin{tabular}{p{0.13\linewidth}p{0.32\linewidth}p{0.38\linewidth}p{0.07\linewidth}}

\toprule
State & HMT coordinate, MeV/\(c^2\) & Subsequent documentary reference & Type\\
\midrule

\(e^\pm\) & \(0.510998950\allowbreak690000809\) & \(0.51099895069(16)\) MeV/\(c^2\) & B\\
\(\mu^\pm\) & \(105.658360294\allowbreak435829303\) & \(105.6583755\pm0.0000023\) MeV/\(c^2\) & F\\
\(\tau^\pm\) & \(1776.860917631\allowbreak432295595\) & \(1776.93\pm0.09\) MeV/\(c^2\) & F\\
\(u,\bar u\) & \(2.165924536\allowbreak173907663\) & \(2.16\pm0.07\) MeV/\(c^2\) & A\\
\(d,\bar d\) & \(4.679550162\allowbreak948874172\) & \(4.70\pm0.07\) MeV/\(c^2\) & A\\
\(c,\bar c\) & \(1270.500838641\allowbreak911135286\) & \(1.2729\pm0.0045\) GeV/\(c^2\) & A\\
\(s,\bar s\) & \(92.940093907\allowbreak845640641\) & \(92.9\pm0.7\) MeV/\(c^2\) & A\\
\(t,\bar t\) & \(172597.479544019\allowbreak136261367\) & \(172.60\pm0.27\) GeV/\(c^2\) & A\\
\(b,\bar b\) & \(4190.119763299\allowbreak790218075\) & \(4.186\pm0.006\) GeV/\(c^2\) & A\\
\(\gamma\) & \(0\) & \(<10^{-18}\) eV/\(c^2\), preserved bound & N\\
\(g\) & \(0\) & Null gauge chart; no experimental comparison row & N\\
\(W^\pm\) & \(80378.964909704\allowbreak547024865\) & \(80.3625\pm0.0077\) GeV/\(c^2\) & F\\
\(Z^0\) & \(91187.533444313\allowbreak407844855\) & \(91.1879\pm0.0020\) GeV/\(c^2\) & F\\
\(H\) & \(125292.466042536\allowbreak133000433\) & \(125.13\pm0.11\) GeV/\(c^2\) & A\\
\(\pi^0\) & \(134.976724327\allowbreak872125739\) & \(134.976827767685\) MeV/\(c^2\), supplement & \(\mathrm F_c\)\\
\(\pi^\pm\) & \(139.570485171\allowbreak572191375\) & \(139.570390983681\) MeV/\(c^2\), supplement & \(\mathrm F_c\)\\
\(K^\pm\) & \(493.677085649\allowbreak530612476\) & \(493.676599458041\) MeV/\(c^2\), supplement & \(\mathrm F_c\)\\
\(K^0\) & \(497.611186497\allowbreak750457359\) & \(497.610767531258\) MeV/\(c^2\), supplement & \(\mathrm F_c\)\\
\(p,\bar p\) & \(938.271751546\allowbreak984898299\) & \(938.27208943\) MeV/\(c^2\), supplement & \(\mathrm F_c\)\\
\(n,\bar n\) & \(939.565810472\allowbreak507023313\) & \(939.56542194\) MeV/\(c^2\), supplement & \(\mathrm F_c\)\\
\end{tabular}
\end{table}
\endgroup

In this table, reference values are presented with the decimal
notation and units declared for each observable. The calculation
record also preserves the decimal string used in the evaluation,
including the difference between that string and its rounded presentation for
the tau, \(Z\) and Higgs. The quark comparison
remains informative until the scheme and scale chart is assembled.
The \(W,Z,H\) sectors also require the resonance convention
to be preserved. The six composite references come from the supplement
identified in the evaluation records; their metadata are not completed
by analogy with the elementary references.

The complete signatures documented in modalities \(\mathrm F\) and
\(\mathrm F_c\), in the order
\(\sigma=(n_A,s_C,k_{\Delta_4},b_{120},b_\zeta)\), are:
\[
\begin{array}{c|r@{,\,}r@{,\,}r@{,\,}r@{,\,}r}
\mu&42&-3&8&0&2\\
\tau&64&0&25&-1&6\\
W&94&-1&0&-2&4\\
Z&95&-1&-11&0&-4\\
\pi^0&44&-4&-25&1&-2\\
\pi^\pm&44&1&-2&2&-1\\
K^\pm&54&-1&17&-2&2\\
K^0&54&0&32&3&-2\\
p&59&0&9&1&0\\
n&59&1&-34&0&1
\end{array}
\]
The angular modalities preserve, respectively for
\(u,d,c,s,t,b,H\), the pairs
\[
(n_A,s_C)=(11,7),(17,8),(61,8),(41,-3),(100,-1),(71,-5),(97,9).
\]
The two declared components are not extended by
three implicit zeros. The evaluation equation and the mechanism
producing a descriptor are different assertions; the proofs of both
are found in their respective definitions and propositions.

\subsection{Relative differences and uncertainty propagation}
\label{subsec:vi-residuales}

For a fixed common chart and two positive coordinates of the same observable,
the relative difference and logarithmic difference are
\[
 \delta_X=\frac{m_X^{\rm HMT}-m_X^{\rm ext}}{m_X^{\rm ext}},
 \qquad r_X=\log\frac{m_X^{\rm HMT}}{m_X^{\rm ext}},
 \qquad \delta_X^{\rm ppm}=10^6\delta_X.
\]
A common unit conversion \(m\mapsto am\), with \(a>0\),
leaves both differences invariant: the factor \(a\) cancels in the
quotients. The absolute difference is multiplied by \(a\), as is
its uncertainty. For a zero external value, a one-sided bound or an
interval without a central estimate, the preceding quotient is replaced by
the comparison appropriate to that observation type.

When a joint uncertainty model is available, let
\(d_X=m_X^{\rm HMT}-m_X^{\rm ext}\). Its variance is
\[
 u^2(d_X)=u^2(m_X^{\rm HMT})+u^2(m_X^{\rm ext})
          -2\operatorname{Cov}(m_X^{\rm HMT},m_X^{\rm ext}).
\]
The formula follows by expanding
\(\operatorname{Var}(U-V)\). If \(u(d_X)>0\), the normalised
difference is \(z_X=d_X/u(d_X)\). Its statistical interpretation requires
the corresponding probabilistic model and covariance; an unrecorded field
is not equivalent to zero variance. Asymmetric uncertainties
are preserved as such and are not automatically converted into a symmetric
deviation.

For several species, the joint covariance matrix is retained. With
the discrete signatures fixed, if \(f_X(\vartheta)=\log(m_X/m_{\rm ref})\)
depends on continuous coordinates already constructed and \(J\) is its
Jacobian matrix, first-order propagation gives
\[
 \Sigma_f=J\Sigma_\vartheta J^{\mathsf T}.
\]
This formula preserves correlations shared by the angular pair,
\(\Delta_4\), the action ratio and the towers. Its application presupposes
the regime in which linearisation controls the remaining terms.
An integer change of signature is recorded as a discrete variation of the
construction, not as an infinitesimal Gaussian fluctuation.

The precision of an evaluation therefore has several sources: the
experimental reference; the unit, scheme and scale chart; the
route distribution; the tower approximation; and the selected
coupling. If \(\Sigma_{ij}\) denotes the covariance between
contributions \(i,j\), the total covariance of their sum is
\(\sum_{i,j}\Sigma_{ij}\). Retaining the cross blocks avoids
ascribing independence to several agreements that use the same
antecedents.

\subsection{Absolute scale and control of the harmonic approximation}
\label{subsec:vi-precision-torres}

The absolute scale and relative correction enter at different
points in the construction. Fix an internally constructed
basal reference section \(m_\star^{(0)}\), and let
\(\chi_\star=\chi_{\rm full}(\sigma_\star,\eta_\star)\).
In the normalisation relative to that section, the expression is
\[
 m_X^\beta=m_\star^{(0)}\,
 \frac{\chi_{\rm full}(\sigma_X,\eta_X)}{\chi_\star}
 R_{\rm act}^{\kappa_X}.
\]
The common basal factor remains explicit. If the reference is electronic,
\(m_\star^{(0)}=m_e^{(0)}\) is the dimensional realisation of
\(\mathfrak e_0\), with all factors in its basal equation, and
\(\chi_\star=\chi_{\rm full}(\sigma_e,\eta_e)\). This recovers
\(m_e^\beta=m_e^{(0)}R_{\rm act}^{\kappa_e}\). If the common
clock is used to write \(m_\star^{(0)}=m_{\rm clk}b_\star\), the
factor \(b_\star\) preserves that basal normalisation; it is not replaced
by \(1\) or incorporated twice into the character. For two states
in the same chart and with the same reference, the quotient is
\[
 \frac{m_Y^\beta}{m_X^\beta}
 =\frac{\chi_{\rm full}(\sigma_Y,\eta_Y)}
        {\chi_{\rm full}(\sigma_X,\eta_X)}
     R_{\rm act}^{\kappa_Y-\kappa_X}.
\]
Cancellation of \(m_\star^{(0)}/\chi_\star\) is exact. It comprises
the common basal prefactor and the common factors in the absolute branch of
action, clock and velocity chart. Consequently,
the action decade \(10^{-34}\) retains its role in
the dimensional realisation and is not inserted again into the relative
corrector. Differences between routes remain in their signatures, readers
and tower sections.

To control a truncation, it is useful to fix the normal form of the character.
Let
\[
 \zeta_{270}=135\Delta_4^2,\qquad
 N_h=\eta_h+b_{120}\mathbf 1_{\{h=120\}},\qquad
 \mathcal S=\langle90,120\rangle\setminus\{0\}.
\]
Then
\[
\begin{split}
 \log\chi_{\rm full}
 ={}&n_AA_{\rm rad}+\frac{s_C}{6}C^*_{\rm rad}
       +k_{\Delta_4}\Delta_4+b_\zeta\zeta_{270}
       +\sum_{h\in\mathcal S}N_hs_h .
\end{split}
\]
In this convention, the height-\(120\) term appears only once.
The term \(b_\zeta\zeta_{270}\) retains its quadratic origin in the
record; \(N_{270}s_{270}\) belongs to the harmonic reading. Equality
of the numerical index does not identify the two objects.

\paragraph{Truncation bound and multiplicative error.}
Suppose that, for \(h>H\), the already determined coefficients of a
route satisfy \(|N_hs_h|\leq Cq^h\), with \(C\geq0\) and
\(0<q<1\). The logarithmic tail \(R_H\) satisfies
\[
 |R_H|\leq\sum_{\substack{h\in\mathcal S\\h>H}}Cq^h
 \leq\sum_{h=H+1}^{\infty}Cq^h
 =\frac{Cq^{H+1}}{1-q}=:\varepsilon_H.
\]
If \(m_H\) is the coordinate obtained by truncating that tower while all other
factors remain fixed, \(m/m_H=e^{R_H}\). Monotonicity of the
exponential yields
\[
 e^{-\varepsilon_H}\leq\frac{m}{m_H}\leq e^{\varepsilon_H},
 \qquad
 \left|\frac{m-m_H}{m_H}\right|\leq e^{\varepsilon_H}-1.
\]
The bound assigns precision to the evaluation from the route and its
operators. The values of \(C,q,H\) are justified by that
construction; the experimental residual selects neither the tower
length nor the bound to be proved.

\subsection{Complete preservation and reproducibility of the comparison}
\label{subsec:vi-catalogo-reproducible}

The catalogue preserves complete records in tabular and
structured formats. Each row retains its source file, ordinal,
original key and occurrence number of that key. This
identification preserves even two literally identical rows:
equality of content does not remove the second occurrence. Format
conversion preserves the original strings, empty cells and
statuses, and includes the hash of the input bytes.

The inventory of \(855\) observables is also linked to the metrological
rows of the integral corpus. The correspondence is verified
by the pair consisting of the observable type and its identifier;
the check compares multiplicities, not merely sets of keys. The internal
tables of cells, collections, multisections and routes remain separate
from that inventory. The \(23\) classes, \(20\) numerical
presentations and historical grid also retain their own files
and statuses.

Transcription checks verify recovery of every
field, row order, duplicates, decimal strings and
correspondence with the reference tables. Their scope is documentary.
The mathematical proofs of the atlas, readers and material realisation
establish the internal outputs; metrological comparison
completes that exposition through well-defined observables and charts.
Documentary integrity and the sufficiency of a proof are
checked separately.

Catalogue extension therefore follows an accumulative rule:
a new evaluation incorporates the state, its record and signature, the tower
section, reader, realisation, comparison chart and
provenance. Earlier cases retain their content and cutoff
date. Thus, inventory growth makes a greater part of the
domain visible without retrospectively altering what each proof and
evaluation established.
